 
\section{The submanifolds of arclength parametrized curves and knots}\label{sec:arclength}
Throughout this section we fix any parameter $s\in (\frac{1}{2},1]$ and
consider the (for $s\in (\frac{1}{2},1)$ fractional) Sobolev space 
\begin{equation}\label{eq:frac-space}
H^{1+s}(\R/\Z,\R^n) \ceq W^{1+s,2}(\R/\Z,\R^n),
\end{equation}
which embeds into the classic H\"older space
$C^{1,s-1/2}(\R/\Z,\R^n)$ and 
hence compactly into $C^1(\R/\Z,\R^n)$; 
see \cite[Theorem A.2]{knappmann-etal_2023}
for a short self-contained proof in the one-dimensional periodic setting
for $s\in (\frac{1}{2},1).$
Therefore, 
the subsets $H^{1+s}_\reg(\R/\Z,\R^n)$ or $H^{1+s}_\inj(\R/\Z,\R^n)$ 
of regular, respectively injective
curves are open subsets of the Hilbert space
$H^{1+s}(\R/\Z,\R^n)$, as well as their
intersection $H^{1+s}_\injreg(\R/\Z,\R^n)$.

%Using a hack to suppress warnings that PDF bookmarks don't like TeX code.
\subsection{The arclength manifolds \texorpdfstring{$\NL^{1+s}$}{N1+s} and \texorpdfstring{$\AL^{1+s}$}{A1+s}}
\label{sec:AL}

To identify the subset 
\begin{equation}\label{eq:arclenth-set}
H^{1+s}_\unitspeed(\R/\Z,\R^n)=\{\g\in H^{1+s}(\R/\Z,\R^n):|\g'|=1\,
\textnormal{
on $\R/\Z$}\}
\end{equation}
of arclength parametrized curves as a submanifold of the open set
$H^{1+s}_\reg(\R/\Z,\R^n)$ we follow the idea of 
\cite[Section 2]{scholtes-etal_2022} and
\cite{knappmann-etal_2023} and introduce
the scalar-valued \emph{logarithmic strain}
$\varSigma:H^{1+s}_\reg(\R/\Z,\R^n)\to H^s(\R/\Z)$ defined as
\begin{equation}\label{eq:logarithmic-strain}
	\varSigma(\g) \ceq \log\big( |\g'|\big)  \in H^s(\R/\Z) \quad\Fo \g\in H^{1+s}_\reg(\R/\Z,\R^n),
\end{equation}
which is a smooth function on the set $H^{1+s}_\reg(\R/\Z,\R^n)$ of
regular curves. It follows from \cite[Proposition 5.1]{knappmann-etal_2023}
for $s\in (\frac{1}{2},1)$
that the differential $\D\varSigma_\g:T_\g H^{1+s}_\reg(\R/\Z,\R^n)
\simeq H^{1+s}(\R/\Z,\R^n)\to T_{\varSigma(\g)}H^s(\R/\Z,\R^n)\simeq
H^s(\R/\Z,\R^n)$ has the form
\begin{equation}\label{eq:log-strain-differential}
	\D\varSigma_\g h=\langle\g',h'\rangle \quad\Fo
	\g\in H^{1+s}_\unitspeed(\R/\Z,\R^n),\, h\in H^{1+s}(\R/\Z,\R^n),
\end{equation}
and is bounded, i.e. has operator norm
\begin{equation}\label{eq:log-strain-bounded-differential}
\|\D\varSigma_\g\|_{\LL(H^{1+s},H^s)}\le C_\varSigma,
\end{equation}
where the constant  $C_\varSigma=C_\varSigma(n,s,\|\g'\|_{H^s})$
depends 
non-decreasingly on the norm $\|\g'\|_{H^s}$
of the unit tangent $\g'$. In view of the explicit
formula \cref{eq:log-strain-differential}
for $\D\varSigma_\g$ it can easily be checked that
the same holds true for $s=1$.
Moreover, \cite[Lemma 5.2]{knappmann-etal_2023}
yields for $s\in (\frac{1}{2},1)$
the existence of a bounded right inverse $Y_\g\in
\LL(H^s(\R/\Z),H^{1+s}(\R/\Z,\R^n))$ of the differential $\D\varSigma_\g$ 
satisfying
$\D\varSigma_\g(Y_\g(f))=f$ for all $f\in H^s(\R/\Z)$,
so that $\D\varSigma_\g$ is surjective for any
$\g\in H^{1+s}_\unitspeed(\R/\Z,\R^n)$. Again, as before, 
the same formula\footnote{Explicitly given in \cref{eq:right-inverse-Y}
in our proof
of \cref{prop:projection-ALs} below.}
for $Y_\g$ given in the proof of \cite[Lemma 5.2]{knappmann-etal_2023}
yields a right inverse of $\D\varSigma_\g$
also in the case $s=1$.
Therefore, we can apply the
preimage theorem\footnote{Notice that the nullspace $\ker(\D\varSigma_\g)$
splits $H^{1+s}(\R/\Z,\R^n)$ because the latter is a Hilbert space.}
(see, e.g., \cite[Theorem 4.J]{zeidler_1993}) 
to
\begin{equation}
\begin{aligned}
	\label{eq:arclength-manifold}
&\NL^{1+s}\ceq H^{1+s}_\unitspeed(\R/\Z,\R^n)=\big(\varSigma|_{H^{1+s}_\reg(\R/\Z,\R^n)}
\big)^{-1}(0),\hspace{5pt}\text{ and }\\
&\AL^{1+s} \ceq H^{1+s}_\injunitspeed(\R/\Z,\R^n)=\big(\varSigma|_{H^{1+s}_\injreg(\R/\Z,\R^n)}\big)^{-1}(0)=\NL^{1+s}\cap H^{1+s}_\inj\left( \R/\Z,\R^n \right).
\end{aligned}
\end{equation}

\begin{theorem}[Arclength manifolds]
\label{thm:arclength-manifold}
Let $s\in (\frac{1}{2},1]$.
The sets $\NL^{1+s}$ of arclength para\-met\-rized curves, and $\AL^{1+s}$ of arclength para\-met\-rized knots are smooth submanifolds
of the open subsets $H^{1+s}_\reg(\R/\Z,\R^n)$ and $ H^{1+s}_\injreg(\R/\Z,\R^n) $, respectively. 
\end{theorem}
To secure compactness of (PS)-sequences  later,
we can fix an arbitrary 
point $\bz\in\R^n$ and
consider the subsets
\begin{equation}\label{eq:arclenth-point}
\NL^{1+s}_\bz\ceq\{\g\in\NL^{1+s}:\g(0)=\bz\},\, \text{ and }\,\AL^{1+s}_\bz\ceq\{\g\in\AL^{1+s}:\g(0)=\bz\}.
\end{equation}
The corresponding ambient space $H^{1+s}_\bz(\R/\Z,\R^n)\ceq\{\g\in
H^{1+s}(\R/\Z,\R^n):\g(0)=\bz\}$ is an affine subspace
and therefore also a smooth submanifold of $H^{1+s}(\R/\Z,\R^n)$, but in order
to come up with a suitable projection later, we express $\NL^{1+s}_\bz$ as
preimage of $(0,\bz)\in H^s(\R/\Z)\times\R^n$, i.e., 
\begin{align*}\label{eq:NL_z-preimage}
	\NL^{1+s}_\bz=\big(
\varSigma|_{H^{1+s}_\reg(\R/\Z,\R^n)}\times P_0\big)^{-1}(0,\bz),
\end{align*}
where the point evaluation $P_0:H^{1+s}_\reg(\R/\Z,\R^n)\to\R^n$ is 
defined by
\begin{equation}\label{eq:point-evaluation}
P_0(\g)\ceq\g(0).
\end{equation}
This linear map is a smooth submersion at every $\g\in H^{1+s}_\reg(\R/\Z,\R^n)$
with differential $\D(P_0)_\g:T_\g H^{1+s}_\reg(\R/\Z,\R^n)\simeq
H^{1+s}(\R/\Z,\R^n)\to T_{P_0(\g)}\R ^n\simeq \R^n$  given by
\begin{equation}\label{eq:P0-differential}
\D(P_0)_\g h=P_0 h=h(0)\quad\Foa h\in H^{1+s}(\R/\Z,\R^n).
\end{equation}
The differential of the product map $\varSigma\times P_0:
H^{1+s}_\reg(\R/\Z,\R^n)\to H^s(\R/\Z)\times\R^n$ is
$$
\D(\varSigma,P_0)_\g=(\D\varSigma_\g,\D(P_0)_\g)
\overset{\cref{eq:P0-differential}}{=}(\D\varSigma_\g,P_0),
$$
 and it possesses as bounded right
inverse the map $Z_\g:H^s(\R/\Z)\times\R^n\to H^{1+s}(\R/\Z,\R^n)$
defined via the right inverse $Y_\g$ of $\D\varSigma_\g$ (see the arguments that lead to
\cref{thm:arclength-manifold}) as
\begin{align*}
	Z_\g(f,\by) \ceq Y_\g(f)+\by-Y_\g(f)(0)\quad\Fo (f,\by)\in
H^s(\R/\Z)\times\R^n.
\end{align*}
Indeed, one computes with the help of \cref{eq:P0-differential}, 
\cref{eq:point-evaluation} and
\cref{eq:log-strain-differential}
\begin{align*}
	\D(\varSigma,P_0)_\g & \big(Z_\g(f,\by)\big)
	\overset{\cref{eq:P0-differential}}{=}
	\big(\D\varSigma_\g
	\big(Y_\g(f)+\by-Y_\g(f)(0)\big),P_0\big(Y_\g(f)+\by-
	Y_\g(f)(0)\big)\big)\\
	&\overset{\cref{eq:log-strain-differential},\cref{eq:point-evaluation}}{=}
	\big(\D\varSigma_\g(Y_\g(f)),\by\big)=(f,\by)\quad\Foa (f,\by)
	\in H^s(\R/\Z)\times\R^n.
\end{align*}
We have thus shown that the differential $\D(\varSigma,P_0)_\g$ is surjective
for any curve $\g\in H^{1+s}_\reg(\R/\Z,\R^n)$, so that we may apply the
preimage theorem as in the proof of \cref{thm:arclength-manifold}.
Including injectivity, we get 
\begin{equation}\label{eq:AL_z-preimage}
	\AL^{1+s}_\bz=\big(
	\varSigma|_{H^{1+s}_\injreg(\R/\Z,\R^n)}\times P_0\big)^{-1}(0,\bz)=\NL^{1+s}_{\bz}\cap H^{1+s}_{\inj}.
\end{equation}

\begin{theorem}[Arclength manifolds with fixed initial point]
\label{thm:arclength-manifold-point}
Let $s\in (\frac{1}{2},1]$ and $\bz\in\R^n$, then the sets $\NL^{1+s}_\bz$ and 
$\AL^{1+s}_\bz$
of arclength parametrized curves and knots $\g$ with fixed initial point $\g(0)=\bz$ are
smooth submanifolds of the open subsets $H^{1+s}_\reg(\R/\Z,\R^n)$ and $ H^{1+s}_{\injreg}(\R/\Z,\R^n) $, respectively. 
\end{theorem}

The tangent space of the arclength manifold $\NL^{1+s}$ at a curve $\g\in
\NL^{1+s}$
is given by
\begin{equation}\label{eq:tangent-space-AL}
	T_\g\NL^{1+s} 
	=
	\ker (\D\varSigma_\g)
	=
	\braces[\big]{ 
		h\in H^{1+s}(\R/\Z,\R^n) : \inner{\g',h'}=0
	},
\end{equation}
whereas for $\NL_\bz^{1+s}$ we have
\begin{equation}\label{eq:tangent-space-ALz}
	T_\g\NL^{1+s}_\bz 
	=
	\ker(\D(\varSigma,P_0)_\g)
	=
	\braces[\big]{ 
		h\in H^{1+s}(\R/\Z,\R^n) : \inner{\g',h'}=0,\,h(0)=0
	}.
\end{equation}
Analogously, one finds $ T_{\g}\AL^{1+s}=T_\g\NL^{1+s} $ for $ \gamma\in \AL^{1+s} $ and $ T_\g\AL^{1+s}_{\bz}=T_\g\NL^{1+s}_{\bz} $ if $ \gamma\in\AL^{1+s}_{\bz} $. 
\begin{proposition}[Bounded linear projection onto the tangent space]
\label{prop:projection-ALs}
Let $s\in (\frac{1}{2},1]$, $ \ML\in \left\{\AL^{1+s},\NL^{1+s}\right\}$, and $\g\in\ML$. Then the mapping
$\pr_\g^{\ML}:H^{1+s}(\R/\Z,\R^n)\to H^{1+s}(\R/\Z,\R^n)$ given by
\begin{equation}\label{eq:projection-ALs}
\pr_\g^{\ML}h\ceq(\Id-Y_\g \D\varSigma_\g)h
\end{equation}
is a linear and bounded projection onto the tangent space $T_\g\ML$
of $\ML$ at $\g$.
\end{proposition}
\begin{proof}
We need the explicit formula for 
the right inverse $Y_\g$ of $\D\varSigma_\g$
as provided in \cite[Lemma 5.2]{knappmann-etal_2023}:
\begin{equation}\label{eq:right-inverse-Y}
	\textstyle
	Y_\g f(t)
	=
	\int_0^t \big( f(u) \, \g'(u) + \P_{\g'(u)}^\perp V_{\gamma}f\big) \d u,
	\;\;
	\text{with}
	\;\;
	V_{\gamma}f=-\varTheta_\g^{-1}\int_0^1 f(u) \, \g'(u)\d u\in\R^n
	,
\end{equation}
where $\varTheta_\g:\R^n\to\R^n$ is 
the bounded invertible linear operator
given by
$\varTheta_\g\ceq\int_0^1\P_{\g'(u)}^\perp \d u$. Note that $ \varTheta_{\gamma}^{-1} $ is indeed bounded because of \cite[Lemma 5.3]{knappmann-etal_2023}, since $ \gamma'\in C^{0,s-\frac{1}{2}}\left( \R/\Z,\R^n \right) $ by Morrey's embedding. 
The vector $V_{\gamma}f$ guarantees that $Y_\g(f)$ is $1$-periodic. 
It depends linearly  on $f$ and satisfies the estimate
\begin{equation}\label{eq:V-h-est}
 |V_{\gamma}f|\le C\|f\|_{C^0}\le C_EC\|f\|_{H^s}\quad\Foa f\in H^s(\R/\Z),
\end{equation}
where the constant $C=C(s,\|\g'\|_{H^s})$ depends non-de\-creas\-ingly on
the norm $\|\g'\|_{H^s}$ and $C_E=C_E(s)$ is the embedding constant from
the compact Morrey embedding $H^s(\R/\Z)\hookrightarrow
 C^0(\R/\Z)$, which is standard for $s=1$; for
 $s\in (\frac{1}{2},1)$ 
 see, e.g., \cite[Theorem A.2]{knappmann-etal_2023}. 
So, we compute
\begin{equation}\label{eq:derivative-right-inverse}
	\pars{Y_\g \, f}'(t) = f(t) \, \g'(t)+\P_{\g'(t)}^\perp V_{\gamma}f;
\end{equation}
$
	\inner{ \pars{Y_\g f}'(t),\g'(t) }=f(t)$ for all $t\in\R/\Z$ and 
$f\in H^s(\R/\Z),
$
and therefore, by \cref{eq:log-strain-differential}
\begin{align}
	\inner{ (\pr_\g^{\ML}h)'(t),\g'(t) }&
	\overset{\cref{eq:log-strain-differential}}{=}
	\inner[\big]{ 
		h'(t) - \pars[\big]{Y_\g \inner{\g', h'}}'(t)
		,
		\g'(t)
	}=0\label{eq:tangent-orthogonal}
\end{align}
for all $h\in H^{1+s}(\R/\Z,\R^n).$
This 
shows that $\pr_\g^{\ML}h\in T_\g\ML$ for all $h\in H^{1+s}(\R/\Z,\R^n).$
Obviously, $\pr_\g^{\ML}$ is linear, so that it is a projection as desired if
we can show that its restriction to $T_\g\ML$ equals the identity (see, e.g., \cite[9.13(1)]{alt_2016}). But linearity of $f\mapsto Y_\g f$ 
indeed yields for $h\in
T_\g\ML$  by means of \cref{eq:tangent-space-AL}
\begin{equation}\label{eq:proj-id}
	\pr_\g^{\ML}h
	=
	h-Y_\g \inner{ \g',h'} 
	\overset{\cref{eq:tangent-space-AL}}{=}h-Y_\g(0)
	=
	h.
\end{equation}
Combining the explicit formulas \cref{eq:right-inverse-Y}
and \cref{eq:derivative-right-inverse} with the estimates
\cref{eq:V-h-est}, \cref{eq:log-strain-bounded-differential} 
and the product rule 
(see, e.g., \cite[Proposition A.5]{knappmann-etal_2023}),     we obtain
\begin{equation}\label{eq:projection-bound}
\|\pr_\g^{\ML}h\|_{H^{1+s}}\le (1+\tilde C C_\varSigma)\|h\|_{H^{1+s}}=:C_\g\|h\|_{H^{1+s}}\quad\Foa 
h\in H^{1+s}(\R/\Z,\R^n),
\end{equation}
where the constant $\tilde C=\tilde C(n,s,\|\g\|_{H^{1+s}})$ and therefore also $C_\g$
depends non-de\-creas\-ingly on
the norm $\|\g\|_{H^{1+s}}.$
\end{proof}
\begin{corollary}[Bounded linear projection onto $T_\g\NL^{1+s}_\bz$]
\label{cor:projection-ALsz}
For $s\in (\frac{1}{2},1]$, $\bz\in\R^n$, $ \ML\in \left\{\NL^{1+s},\AL^{1+s}\right\} $, and $\g\in\ML_\bz\ceq\{\eta\in \ML: \eta(0)=\bz\}$, the
mapping $\pr_\g^{\ML_\bz}:H^{1+s}(\R/\Z,\R^n)\to 
H^{1+s}(\R/\Z,\R^n)$
defined by
\begin{equation}\label{eq:projection-ALsz}
\pr_\g^{\ML_\bz}h\ceq\pr_\g^{\ML}h-
\pr_\g^{\ML}h(0)
\end{equation}
is a linear and bounded projection onto the tangent space
$T_\g\ML_\bz$ of $\ML_\bz$ at $\g$. 
Here, $\pr_\g^{\ML}$ is the projection provided by \cref{prop:projection-ALs}.
\end{corollary}
\begin{proof}
Notice that the
linearity and boundedness of $\pr_\g^{\ML}$ 
transfers to $\pr_\g^{\ML_\bz}$, and with 
$(\pr_\g^{\ML_\bz}h)'
=
(\pr_\g^{\ML}h)'$ and $\pr_\g^{\ML_\bz}h(0)=0$ we find
by \cref{eq:tangent-orthogonal} and
\cref{eq:tangent-space-ALz} that $\pr_\g^{\ML_\bz}h
\in T_\g\ML_\bz$. Finally, for $h\in T_\g\ML_\bz
\subset T_\g\ML$ we have $h(0)=0$ and therefore, due to \cref{eq:proj-id}, 
\begin{align*}
	\pr_\g^{\ML_\bz}h=
	\pr_\g^{\ML}h\overset{\cref{eq:proj-id}}{=}h.
\end{align*}
\end{proof}

Now \cref{thm:arclength-manifold,thm:arclength-manifold-point},
\cref{prop:projection-ALs}, and \cref{cor:projection-ALsz} 
together with the estimate \cref{eq:projection-bound}
verify the basic abstract conditions
\ref{condE} and \ref{condP} for the arclength manifolds $\AL^{1+s}$ 
and
$\AL^{1+s}_\bz$.
\begin{corollary}[\ref{condE} and \ref{condP} for arclength manifolds]
\label{cor:basic-assumptions}
Let $s\in (\frac{1}{2},1]$ and $\bz\in\R^n$ be fixed. Then condition \ref{condE}
holds true for $\ML\in \left\{\AL^{1+s},\AL_{\bz}^{1+s},\NL^{1+s},\NL_{\bz}^{1+s}\right\}$, if we choose 
$\OL \ceq H^{1+s}_\reg(\R/\Z,\R^n)$, or $ \OL\ceq H^{1+s}_\injreg(\R/\Z,\R^n)$ respectively, as an open subset of the Hilbert space
$\HL \ceq H^{1+s}(\R/\Z,\R^n)$, which embeds compactly into
the Banach space $\BL \ceq C^1(\R/\Z,\R^n)$. 
Moreover, the linear projections $\pr_\g^{\ML}:
T_\g\OL\to T_\g\ML$ for $\g\in\ML$ have property \ref{condP}.
\end{corollary}

\begin{remark}[Explicit Lagrange multiplier]\label{rem:explicit-lagrange-multiplier}
	The concrete formulas for the right inverse $ Y_{\gamma} $ in \cref{eq:right-inverse-Y} and for the projection $ \pr^{\ML}_{\gamma} $ in \cref{eq:projection-ALs} for $ \ML\in \{\AL^{1+s},\NL^{1+s}\} $ allow us to derive an explicit expression for the Lagrange multiplier arising from the arclength constraint. This turns out to be extremely useful in the regularity theory in \cref{sec:Smoothness}. 

	Suppose $ \gamma\in \ML $ is critical for the restricted energy $ \widetilde{\EL}\ceq\EL\vert_{\ML} $, where $ \EL \in C^1(\OL)$ for an open subset of the Hilbert space $ \HL=H^{1+s}\left( \R/\Z,\R^n \right) $ with $ \ML\subset\OL $. Then $ \D \EL_{\gamma}h=\D \widetilde{\EL}_{\gamma}h =0$ for all $h\in \T_{\gamma}\ML $, which is equivalent to 
	\begin{equation}\label{eq:Lagra_Mult}
	\begin{aligned}
		0=\D\EL_{\gamma} \pr^{\ML}_{\gamma}\overset{\cref{eq:projection-ALs}}{=} \D\EL_{\gamma}\left( \Id_{\HL}-Y_{\gamma}\D\Sigma_{\gamma} \right).
	\end{aligned}
	\end{equation}
	By the Lagrange multiplier theorem, see \cite[Corollary 3.5.29]{abraham-etal_1988}, there exists $ \Lambda\in \left( H^{s}(\R/\Z) \right)^* $ such that $ \D\EL_{\gamma}-\Lambda\D\Sigma_{\gamma}=0 $. This in combination with \cref{eq:Lagra_Mult} yields the formula 
	\[\Lambda=\D\EL_{\gamma}Y_{\gamma}.\]
\end{remark}

\subsection{The topology of arclength manifolds}
\label{sec:AL-toplogy}
Fix a point $\bz\in\R^n$ and $s\in \intervaloc{0,1}$. The Hilbert space $ \HL=H^{1+s}\left( \R/\Z,\R^n \right) $ with its inner product\footnote{For $ s=1 $ simply replace the double integral in \cref{eq:fractionalscalarproduct} by $ \left\langle \gamma'',h'' \right\rangle_{L^2}  $.}
\begin{equation}\label{eq:fractionalscalarproduct}
\begin{aligned}
	\inner{f,h}_\HL&\ceq\inner{f,h}_{H^{1+s}}=\inner{f,g}_{L^2}+\inner{f',g'}_{L^2}\\
	&+\int_{\R/\Z}\int_{\R/\Z}\frac{\inner{f(u)-f(v),h(u)-h(v)}}{
	\abs{u-v}_{\R/\Z}^{1+2s}}\d u \d v,\quad 
	\text{for $f,h \in \HL$},
\end{aligned}
\end{equation}

induces a Riemannian metric $g$ on  the arclength submanifolds
\begin{align*}
	&\NL^{1+s}_\bz\subset\NL^{1+s}\subset H^{1+s}_\reg(\R/\Z,\R^n)\subset
H^{1+s}(\R/\Z,\R^n)\,\text{ and }\\
&\AL^{1+s}_\bz\subset\AL^{1+s}\subset H^{1+s}_\injreg(\R/\Z,\R^n)\subset
H^{1+s}(\R/\Z,\R^n).
\end{align*}
More precisely, let $ \ML\in \left\{\NL^{1+s},\NL^{1+s}_{\bz},\AL^{1+s},\AL^{1+s}_{\bz}\right\} $,
$g_\g(f,h)\ceq\langle f,h\rangle_\HL$
for $\g\in\ML$ and $f,h\in T_\g\ML$. So, $\ML$ is a
smooth Riemannian submanifold of the open set $H^{1+s}_\reg(\R/\Z,\R^n) $, or $H^{1+s}_\injreg(\R/\Z,\R^n)$,
respectively; see
\cite[Theorem 3.6]{palais_1966}.
As in the abstract setting of \cref{sec:grad_flows} we have
the path metric (cf.{} \cref{eq:path-metric})
\begin{equation}\label{eq:path-metricALsz}
\dist_{\ML}(\g,\eta)\ceq\inf_{c\in \PS(\g,\eta)}\int_0^1 
\sqrt{g_{c(t)}(\dot{c}(t),\dot{c}(t))}\d t=
\inf_{c\in \PS(\g,\eta)}\int_0^1 \|\dot{c}(t)\|_\HL \d t 
\end{equation}
measuring the distance of
two curves $\g,\eta\in\ML$
in the same connected component of $\ML$, where 
\begin{equation}\label{eq:path-spaceAL}
	\PS(\g,\eta)
	\ceq
	\braces[\big]{
		c\in C^1([0,1],\ML)
		:
		c(0)=\g,c(1)=\eta
	}.
\end{equation}

To investigate the topology of $\ML$ 
more closely,
we look at a mapping that transforms an arbitrary regular closed curve $\g
\in H^{1+s}_\reg(\R/\Z,\R^n)$
into an arclength parametrized curve. The arclength $\psi_\g:
[0,1]\to [0,1]$
(rescaled to the unit interval and then extended to all of $\R$ such that 
$\psi_\g(t+1)=1+\psi_\g(t)$ for all $t\in\R$)  
is given by the strictly monotonic function 
$\psi_\g(t)\ceq\frac1{\LL(\g)}\LL(\g|_{[0,t]})$ for $t\in [0,1]$.
It has a strictly monotonic inverse
$\psi_\g^{-1}:[0,1]
\to [0,1]$.
We prove in \Cref{thm:arclength} 
in \cref{sec:arclength-reparam}
that the reparametrization and rescaling mapping
$\ARC:H^{1+s}_\reg(\R/\Z,\R^n)\to H^{1+s}_\unitspeed(\R/\Z,\R^n)$
given by
\begin{equation}\label{eq:ARC}
	\ARC(\g)
	\ceq
	\sdfrac{1}{\LL(\g)}\,\g\circ \psi_\g^{-1}\quad\Fo \g\in H^{1+s}_\reg(\R/\Z,\R^n)
\end{equation}
is continuous with respect to the $H^{1+s}$-topology for any $s\in (\frac{1}{2},1)$.
For $s=1$ this was proven by Reiter in \cite[Theorem 1.13]{reiter_2009}.
Notice that rescaling and reparametrization does not destroy the injectivity
of a curve so that we have the inclusion
$\ARC(H^{1+s}_\injreg(\R/\Z,\R^n))\subset
\AL^{1+s}$ for all $s\in (\frac{1}{2},1]$.

\begin{theorem}[Arclength manifolds are strong deformation retracts]
\label{thm:strong-deformation-ALs}
The open set\newline $H_\reg^{1+s}(\R/\Z,\R^n)\subset H^{1+s}(\R/\Z,\R^n)$ strongly
deformation retracts onto the submanifolds $\NL^{1+s}$ and 
$\NL^{1+s}_\bz$
for any $s\in (\frac{1}{2},1]$ and $\bz\in\R^n$, whereas $ H^{1+s}_{\injreg}\left( \R/\Z,\R^n \right) $ strongly deformation retracts onto $ \AL^{1+s} $ and $ \AL^{1+s}_{\bz} $. 
\end{theorem}

\begin{proof}
If $\g\in\NL^{1+s}$ then $\LL(\g)=1$ and $\psi_\g(t)=t$ for any $t\in\R$ so
that $\ARC(\g)=\g$ by definition \cref{eq:ARC}. 
Therefore, the mapping $\ARC \colon H^{1+s}_\reg(\R/\Z,\R^n)\to
H^{1+s}_\unitspeed(\R/\Z,\R^n)=\NL^{1+s}$ 
is continuous with $\ARC|_{\NL^{1+s}}=
\Id_{\NL^{1+s}}$, whence a retraction onto $\NL^{1+s}$.
For a given point $\bz\in\R^n$ we simply compose
the map $\ARC$ with the smooth
shift $\tau_\bz:H^{1+s}(\R/\Z,\R^n)\to
H^{1+s}(\R/\Z,\R^n)$ defined as $\tau_\bz(h) \ceq h+\bz-h(0)$
satisfying $\tau_\bz(\g)=\g$ for all $\g\in\AL^{1+s}_\bz$ 
so that we
obtain the retraction $\tau_\bz\circ\ARC:H^{1+s}_\injreg(\R/\Z,\R^n)
\to\NL^{1+s}_\bz$.

Now, we need to construct a homotopy between the identity map on the set
$\OL \ceq H^{1+s}_\reg(\R/\Z,\R^n)$ and the retraction onto $\NL^{1+s}$, and $\NL^{1+s}_\bz$, respectively.
Consider for any curve $\eta\in \OL$
the family of parameter transformations 
\begin{equation}\label{eq:familiy-param}
\phi_\eta(\sigma) \ceq (1-\sigma)\Id_\R+
\sigma\psi_\eta^{-1}:[0,1]\to [0,1]\quad\Fo \sigma\in [0,1],
\end{equation}
where, as before, $\psi_\eta^{-1}$ is the strictly monotonic inverse
of the rescaled arclength $\psi_\eta(t)=\frac1{\LL(\eta)}\LL(\eta|_{[0,t]})$.
As before, we may extend $\phi_\eta$ to all of $\R$ with
$\phi_\eta(t+1)=1+\phi_\eta(t)$ for all $t\in\R$.
Clearly, 
$\phi_\eta(\sigma)(0)=0$ and $\phi_\eta(\sigma)(1)=1$, and $\phi_\eta(\sigma)
(\cdot)$
is strictly
increasing for all $\sigma\in [0,1]$, since
$\phi_{\eta}(\sigma)'(t)=(1-\sigma)+\sigma (\psi_\g^{-1})'(t)>0$ for
any $t\in [0,1]$. We can interpret this family $\phi_\eta(\sigma)$ as 
a mapping $\phi_{\cdot}(\cdot):\OL\times [0,1]\to
H^{1+s}_\reg([0,1])$, which clearly is continuous on
$\OL\times [0,1]$ 
due to \Cref{cor:restriction-inverse-operator} and 
\Cref{lem:parametrization_continuous} in \cref{sec:arclength-reparam}.
Next, we modify the mapping $\ARC$ in \cref{eq:ARC} to obtain
a family of mappings $\CAPARC(\eta,\sigma)$ for $\sigma\in [0,1]$ defined as
\begin{equation}\label{eq:ARC-sigma}
\textstyle
\CAPARC(\eta,\sigma)\ceq\frac1{1+\sigma(\LL(\eta)-1)}\,\eta\circ\phi_\eta(\sigma)
\quad\Fo\eta\in \OL.
\end{equation}
\Cref{composition_continuous} in \cref{sec:arclength-reparam}
implies that the map 
$\CAPARC(\cdot,\cdot):\OL\times [0,1]\to \OL
$ is continuous. With $\phi_\eta(0)(t)=t$ for all $t\in\R$ one has $\CAPARC(\eta,0)=\eta$
for all $\eta\in \OL$, and $\CAPARC(\eta,1)\in\NL^{1+s}$ since $\phi_\eta(1)=
\psi_\eta^{-1}$ so that $\CAPARC(\eta,1)=\ARC(\eta)$ by \cref{eq:ARC}, 
and the original map
$\ARC$ was shown 
to be a retraction onto $\NL^{1+s}$.
For $\g\in\NL^{1+s}$ one has $\LL(\g)=1$ and $\psi_\g^{-1}(t)=t$ whence 
$\phi_\g(\sigma)(t)=t$ for all $t\in\R$ and $\sigma\in [0,1]$, 
and therefore 
$
\CAPARC(\g,\sigma)=\frac1{1+\sigma(1 -1)}\,\g=\g.$
Hence, $\CAPARC(\cdot,\cdot)$ is a strong deformation 
retraction of $\OL$ onto $\NL^{1+s}$.

If we fix a point $\bz\in\R^n$ to treat the arclength submanifold
$\NL^{1+s}_{\bz}$ instead of $\NL^{1+s}$, then we take 
the family of shifts $T_\bz(h,\sigma)\ceq
h+\sigma(\bz-h(0))$ for $h\in\OL$. 
Then the composition
$T_\bz(\CAPARC(\cdot,\sigma),\sigma)$ for $\sigma\in [0,1]$
yields a strong deformation retraction of $\OL$ onto $\NL^{1+s}_\bz$.
Indeed, we have
$T_\bz(h,0)=h$, 
$T_\bz(h,1)=h+\bz-h(0)$ hence $T_\bz(h,1)'=
h'$, and
$T_\bz(h,1)(0)=\bz$ for all $h\in \OL$.
So, $T_\bz(\CAPARC(\eta,0),0)=\CAPARC(\eta,0)=\eta$, and
$T_\bz(\CAPARC(\eta,1),1)=\CAPARC(\eta,1)+\bz-\CAPARC(\eta,1)(0)
\in\NL^{1+s}_\bz$ for every $\eta\in\OL$. In addition,
for $\g\in\NL^{1+s}_\bz$ we have shown above
that $T_\bz(\CAPARC(\g,\sigma),\sigma)=
T_\bz(\g,\sigma),$
which equals $\g$
since $\g(0)=\bz$.

Replacing $ \NL^{1+s} $ and $ H^{1+s}_{\reg} $ by $ \AL^{1+s} $ and $ H^{1+s}_{\injreg} $ respectively, yields the desired statement for the submanifolds of arclength parametrized knots. 
\end{proof}
\begin{remark}\label{rem:topological-justification-AL}
\cref{thm:strong-deformation-ALs} 
justifies restricting knot energies to the submanifold of arclength parametrized
curves topologically: As the knot space $H^{1+s}_\injreg(\R/\Z,\R^n)$ strongly
deformation retracts onto $\AL^{1+s}$ and also onto $\AL^{1+s}_\bz$, all
three topological spaces are homotopy equivalent and hence contain
the same topological information. 
\end{remark}
The last result of this section gives a justification from the perspective of knot theory. 
\begin{theorem}[Connected components are knot classes]
\label{thm:components-knot-classes}
Let $s\in (\frac{1}{2},1]$. If two knots $\g_0$ and $\g_1$ lie in the same
connected component of $\AL^{1+s}$ then they are $C^1$-ambiently isotopic.
If, on the other hand,
two knots $\g_0,\g_1\in\AL^{1+s}$ are $C^1$-isotopic in 
$\R^n$,
then they lie in the same connected component
of $\AL^{1+s}$.
All statements remain 
true if we replace $\AL^{1+s}$ by $\AL^{1+s}_\bz$ for some fixed point
$\bz\in\R^n$.
\end{theorem}
\begin{proof}
Fix $s\in (\frac{1}{2},1]$.
It suffices to treat $\AL \ceq \AL^{1+s}$.
If $\g_0,\g_1\in\AL$ are contained in the same connected component
of the smooth submanifold $\AL\subset\OL\ceq
H^{1+s}_\injreg(\R/\Z,\R^n)$ then there exists a path $c\in  
C^\infty([0,1],\AL)$ such that
$c(0)=\g_0$ and $c(1)=\g_1$; see, e.g., \cite[Lemma 3.1]{palais_1966}.
Set $J(t,\sigma) \ceq c(\sigma)(t)$ for $t\in\R/\Z$ and $\sigma\in [0,1]$, to
find that the curve
$J(\cdot,\sigma)\in\AL$ is a $C^1$-embedding for any $\sigma\in [0,1]$
by means of the Morrey embedding $H^{1+s}(\R/\Z,\R^n)\hookrightarrow
C^1(\R/\Z,\R^n)$; see \cite[Theorem A.2]{knappmann-etal_2023}. 
Moreover, one can show with this embedding 
that $J$ is also of class $C^1(\R/\Z\times [0,1],
\R^n)$ so that $J$ constitutes a $C^1$-isotopy with $J(\cdot,0)=\g_0$ and
$J(\cdot,1)=\g_1$. According to \cite[Theorem 1.2]{blatt_2009a} the two knots
$\g_0$ and $\g_1$ are therefore ambiently isotopic, i.e., in the same knot
class.

Let $J \in C^1(\Circle\times[0,1] ,\AmbSpace)$ be an isotopy satisfying $J(\cdot,0) = \gamma_0$ and $J(\cdot,1) = \gamma_1$. 
We consider the metric space $X \ceq C^1(\Circle,\AmbSpace)$ with the $C^1$-norm.
Define $\varGamma \in C^0([0,1],X)$ by $\varGamma(t) \ceq J(\cdot,t)$.
The set $U \ceq C^1_{\mathrm{ir}}(\Circle,\AmbSpace)$ is open in $X$,
and $\varGamma$ maps into it.
Observe that the distance function
\[
    \dist(\gamma, X \setminus U) \ceq \smash{\inf_{\eta \in X \setminus U} \|\gamma - \eta\|_{C^1}} \quad \text{ for }\gamma\in U
\]
is Lipschitz continuous.
Thus, the function $F \colon \intervalcc{0,1} \to \R$, $F(t) \ceq \dist(\varGamma(t), X \setminus U)$ is continuous. 
The interval $\intervalcc{0,1}$ is compact, so $F$ must have a minimizer $t_0 \in \intervalcc{0,1}$.
Because of $\varGamma(t_0) \in U$ and because $U$ is open, the minimum value must be positive: $2\varepsilon\ceq F(t_0) > 0$. 
Since $C^\infty(\Circle\times[0,1],\AmbSpace)$ is dense in $C^1(\Circle\times[0,1],\AmbSpace)$, there is an  $H \in C^\infty(\Circle\times[0,1],\AmbSpace)$ so that
$
    \norm{J - H}_{C^1(\Circle\times[0,1],\R^n)}
    <
    \varepsilon
$.
For the path $G$ defined by $G(t) \ceq H(t,\cdot)$ this means that
\[
\norm[\big]{ 
		t \mapsto \norm{ \varGamma(t) - G(t) }_{C^1} 
	}_{L^\infty} 
    \leq 
    \norm{J - H}_{C^1([0,1] \times \Circle,\R^n)}
    <
    \varepsilon
	.
\]
By construction, we have for every $t \in [0,1]$ that
\[
    G(t) 
	\in 
	B_\varepsilon(\varGamma(t)) \cap C^2(\Circle,\AmbSpace)
	\subset
	U \cap C^2(\Circle,\AmbSpace) 
    = C^2_{\mathrm{ir}}(\Circle,\AmbSpace)
	.
\]
Since the embedding $C^2(\Circle,\AmbSpace) \hookrightarrow H^{1+s}(\Circle,\AmbSpace)$ is linear and continuous, the path $G$ is an element of
$C^\infty([0,1],H^{1+s}_{\mathrm{ir}}(\Circle,\AmbSpace))$.
Now we define the linear homotopies
\[
    \varGamma_1(t) \ceq (1-t) \, \gamma_0 + t \, G(0)
    \quad \text{and} \quad 
    \varGamma_2(t) \ceq (1-t) \, G(1) + t \, \gamma_1
    ,
    \quad 
    \text{for $t \in [0,1]$}
    .
\]
Notice that $G(0) \in B_\varepsilon(\gamma_0)$, $G(1) \in B_\varepsilon(\gamma_1)$, that these balls contained in $ C^1_{\injreg}(\Circle,\AmbSpace) $ are convex,  and that knot classes are stable under small $ C^1 $-pertubations, see \cite{blatt_2009a}. 
Thus, we have $\varGamma_1, \varGamma_2 \in C^\infty([0,1],H^{1+s}_{\mathrm{ir}}(\Circle,\AmbSpace))$.
Concatenation of the paths $\varGamma_1$, $G$, and $\varGamma_2$ yields a continuous path from $\gamma_0$ to $\gamma_1$ in $H^{1+s}_{\mathrm{ir}}(\Circle,\AmbSpace)$.
Finally, we retract this path to a continuous path in $\mathcal{A}$, by means of \cref{thm:strong-deformation-ALs}.
\end{proof}

\subsection{Critical knots in arclength submanifolds}\label{sec:AL-critical}
In this abstract setting we consider for $s\in (\frac{1}{2},1]$  
an arbitrary  continuously differentiable parameter-invariant
energy 
$$
\EL \colon \OL \ceq H^{1+s}_\reg(\R/\Z,\R^n)\to\R.
$$
The latter means that 
$\EL(\eta)=\EL(\eta\circ\psi)$ for any $\eta\in\OL$,
where $\psi\in H^{1+s}([0,1])$ is a
strictly increasing parameter transformation with $\psi(0)=0$ and $\psi(1)=1$,
which can be extended to all of $\R$ to satisfy $\psi(t+1)=1+\psi(t)$
for all $t\in\R$.  

One of the simplest examples of such a parameter-invariant functional is the
length $\LL \colon \OL\to\R$, which is actually a smooth submersion, since its differential
$\D\LL_\eta\colon T_\eta\OL\simeq H^{1+s}(\R/\Z,\R^n)\to T_{\LL(\eta)}\R\simeq\R$
satisfies $\D\LL_\eta(\eta)=\LL(\eta)$ for all $\eta\in\OL$ because of the 
positive
$1$-homogeneity of $\LL$ and Euler's homogeneous function 
theorem. 
Hence, by the preimage theorem \cite[Theorem 4.J]{zeidler_1993}, the set
\begin{equation}\label{eq:Omega-s}
\varOmega^{1+s}\ceq\{\g\in\OL:\LL(\g)=1\}=\LL^{-1}(1),
\end{equation}
is a smooth submanifold of $\OL$
with its tangent space given by
\begin{equation}\label{eq:TgOmega-s}
	T_\g\varOmega^{1+s}
	=
	\bigg\{
		\nu\in H^{1+s}(\R/\Z,\R^n) 
		: 
		\D\LL_\g(\nu)=\int_0^1\sdfrac{\langle\g'(u),\nu'(u)\rangle}{|\g'(u)|}\d u=0
	\bigg\}.
\end{equation}
We also consider the smooth submanifolds with a fixed point $\bz\in\R^n$:
\begin{equation}\label{eq:Omega-sz}
\varOmega^{1+s}_\bz\ceq\{\g\in\varOmega^{1+s}:\g(0)=\bz\}
\end{equation}
with tangent space
\begin{equation}\label{eq:TgOmega-sz}
T_\g\varOmega^{1+s}_\bz=\{\nu\in T_\g\varOmega^{1+s}:\nu(0)=0\}.
\end{equation}
We now relate the arclength manifolds $\NL^{1+s}$ and $\NL^{1+s}_{\bz}$ 
to $\varOmega^{1+s}$ and
$\varOmega^{1+s}_{\bz}$ via general parameter-invariant energies $\EL\in C^1(\OL)$
and their critical points under a fixed-length constraint.

Note that as before, since $ \AL^{1+s} $ and $ \AL^{1+s}_{\bz} $ are open subsets of $ \NL^{1+s} $ and $ \NL_{\bz}^{1+s} $ respectively, the results of this section hold as well for arclength parametrized knots if we choose $ \OL$ to be the open set $ H^{1+s}_{\injreg}(\R/\Z,\R^n) $. This is due to the fact that the corresponding tangent spaces coincide; see \cref{eq:tangent-space-AL,eq:tangent-space-ALz}.
Only $|\g'|=1$ will be used in the arguments of this subsection, injectivity of $\g$ does not play a r{\^o}le here.

\begin{theorem}[Arclength criticality vs. criticality under length constraint]
\label{thm:critical-points}
Let $s\in (\frac{1}{2},1]$ and $ \ML\in \{\NL^{1+s},\AL^{1+s}\} $.
If $\g\in\ML\cap H^{2+s}(\R/\Z,\R^n)$ is a critical point of the restricted
parameter-invariant energy $\widetilde{\EL}\ceq\EL|_{\ML}$, 
then $\g$ is a critical
point of $\EL-\lambda\LL$ on $\OL$ for some $\lambda\in\R$. If  
$\EL$ is 
positively $\alpha$-homogeneous for some 
$\alpha\in\R$, then
$\lambda = \alpha \, \EL(\g) = \alpha \, \widetilde{\EL}(\g)$. In particular, if $\alpha\not=0$ and
$\widetilde{\EL}(\g)\not= 0$ then $\lambda\not= 0$, whereas
if $\EL$ is scale-invariant or $\widetilde{\EL}(\g)=0$, then $\lambda=0$.
\end{theorem}

Notice that we require higher regularity of the critical point, 
which we will verify for the energies in
our applications in geometric knot theory in \cref{sec:Smoothness}. 
This higher regularity  is necessary in
the following decomposition of tangent vectors $\nu\in T_\g\varOmega^{1+s}$
into a tangential and a non-tangential part  with respect to $\g$. 
We lose one derivative of $\g$ when we split off the tangential part.
\begin{lemma}[Decomposition of tangent vectors]
\label{lem:decomposition}
For any $\g\in \NL^{1+s}\cap H^{2+s}(\R/\Z,\R^n)\subset\varOmega^{1+s}$ 
and $\nu\in T_\g\varOmega^{1+s}$ there exists
$h\in T_\g\NL^{1+s}$ and a function $\varrho\in H^{1+s}(\R/\Z)$ with $\varrho(0)=\varrho(1)=0$, such 
that 
\begin{equation}\label{eq:decomposition}
	\nu(t)=h(t) + \varrho(t) \, \g'(t)\quad\Foa t\in\R/\Z.
\end{equation}
Likewise, if $\g\in\NL^{1+s}_{\bz}\cap H^{2+s}(\R/\Z,\R^n)$ for some fixed $\bz\in\R^n$
then every $\nu\in T_\g\varOmega^{1+s}_\bz$ can be decomposed as in \cref{eq:decomposition}
where now $h\in T_\g\NL^{1+s}_\bz$.
\end{lemma}
\begin{proof}
Set $\varrho(t)\ceq\int_0^t \inner{\nu'(u),\g'(u)}_{\R^n} \dd u$, which is of class
$H^{1+s}(\R/\Z)$ due to the product rule (see, e.g., \cite[Proposition A.5]{knappmann-etal_2023}), and satisfies $\varrho(0)=0=\varrho(1)$ since $\nu\in T_\g\varOmega^{1+s}$ and
$\abs{\g'}=1$ on $\R/\Z$. The function $h\ceq\nu-\varrho\g'$ is in $H^{1+s}(\R/\Z,\R^n)$
because of our regularity assumption on $\g$ and by the product rule, and one
computes (abbreviating $\inner{\cdot,\cdot} \ceq \inner{\cdot,\cdot}_{\R^n}$)
\begin{align*}
	\inner{h',\g'} 
	&=
	\inner{\nu',\g'} -\varrho' \inner{\g',\g'}-\rho \inner{\g'',\g'}
	=
	\inner{\nu',\g'} - \inner{\nu',\g'} \inner{\g',\g'}
	=
	0,
\end{align*}
where we used the definition of $\varrho$ and the fact that $\g$ is arclength parametrized.
Therefore, $h\in T_\g\NL^{1+s}$ (see \cref{eq:tangent-space-AL}), which establishes the decomposition \cref{eq:decomposition}.

With the same Ansatz for $\varrho$ one obtains for $\g\in\NL^{1+s}_\bz$ 
and $\nu
\in T_\g\varOmega_\bz^{1+s}$ that $h(0)=\nu(0)-\varrho(0)\g'(0)=0$ so that 
$h\in T_\g\NL^{1+s}_\bz$ by means of \cref{eq:tangent-space-ALz}.
\end{proof}
With this decomposition of tangent vectors of $\varOmega^{1+s}$ and 
$\varOmega^{1+s}_\bz$ we can
now prove criticality under the fixed-length constraint.

\begin{refproof}[Proof of \cref{thm:critical-points}]
As mentioned earlier, it suffices to prove the theorem for $ \ML=\NL^{1+s} $. For arbitrary $\nu\in T_\g\varOmega^{1+s}$ we find 
according to \cref{lem:decomposition}
some $h\in T_\g\ML$ and a function $\varrho\in H^{1+s}(\R/\Z)$ with $\varrho(0)=\varrho(1)=0$ such that
\cref{eq:decomposition} holds true. Consider the family $\varphi_\tau\ceq
\Id_\R + \tau\varrho\in H^{1+s}([0,1])$ satisfying $\varphi_\tau(0)=0$ and
$\varphi_\tau(1)=1$, and 
$$
\varphi_\tau'(t)=1+\tau\varrho'(t)\ge 1 - \abs{\tau} \norm{\varrho'}_{C^0}\ge
1-\tau_0 \,C_E \norm{\varrho}_{H^{1+s}}>0\quad\Foa t\in [0,1],
$$
if $|\tau|\le\tau_0\ll 1$. Here, $C_E$ denotes the constant of the Morrey embedding
theorem \cite[Theorem A.2]{knappmann-etal_2023}. So, $\psi\ceq\varphi_\tau$ is an
admissible strictly increasing parameter transformation under which $\EL$ is assumed to be invariant, i.e.,
the function $\tau\mapsto \EL(\g\circ\varphi_\tau)$ is constant on $(-\tau_0,\tau_0)$.
With 
the additional properties 
$$
	\textstyle
	\frac{\partial}{\partial \tau}\varphi_\tau(t) |_{\tau=0}
	=
	\varrho(t)
	\quad\AND\quad
	\varphi_0(t) = t \quad\Foa t\in [0,1],
$$
we conclude
\begin{align*}
	0 &
	\textstyle
	= 
	\frac{\partial}{\partial \tau}\EL(\g\circ\varphi_\tau) \big|_{\tau=0}
	= 
	\D\EL_\g \!\pars[\big]{ \pars{\g'\circ\varphi_0} \cdot\frac{\partial}{\partial \tau}\varphi_\tau |_{\tau=0} }
	=
	\D\EL_\g(\varrho\g')
	\\
	&\overset{\cref{eq:decomposition}}{=}
	\D\EL_\g(\nu-h)
	=
	\D\EL_\g(\nu)-\D\EL_\g(h)
	=
	\D\EL_\g(\nu)-\D\widetilde{\EL}_\g(h)
	=
	\D\EL_\g(\nu)
	.
\end{align*}
Since $\nu\in T_\g\varOmega^{1+s}$ 
was arbitrary we infer from the Lagrange multiplier 
theorem (see \cite[Corollary 3.5.29]{abraham-etal_1988}) 
that there is some number $\lambda\in\R$ such that
$(\D\EL-\lambda \D\LL)_\g=0$.

For the evaluation of $\lambda$ we use that $\LL(\g)=1$,
the positive
$1$-homogeneity of $\LL$, and the positive $\alpha$-homogeneity of $\EL$ to
compute with Euler's homogeneous function theorem
$
\lambda = \lambda \, \LL(\g) = \lambda \D\LL_\g(\g) = \D\EL_\g(\g) = \alpha \, \EL(\g) =
 \alpha \, \widetilde{\EL}(\g).
$
\end{refproof}

Fixing a point $\bz\in\R^n$ adds a vectorial constant Lagrange multiplier.
\begin{corollary}[Criticality with fixed point]
\label{cor:critical-points}
Let $s\in (\frac{1}{2},1]$, $\bz\in\R^n$ and $ \ML_{\bz}\in \{\NL^{1+s}_{\bz},\AL^{1+s}_{\bz}\} $. If $\g\in\ML_{\bz}\cap
H^{2+s}(\R/\Z,\R^n)$ is a critical knot of the restricted 
parameter-invariant energy
$\widetilde{\EL}\ceq\EL|_{\ML_{\bz}}$, then $\g$ is critical for
$\EL-\lambda \, \LL- \inner{\Upsilon,P_0(\cdot)}$ for some
$\lambda\in\R$ and $\Upsilon\in\R^n$. 
If $\EL$ is 
positively $\alpha$-homogeneous for some $\alpha\in\R$,
then $\lambda + \inner{\Upsilon,\bz}=\alpha \, \widetilde{\EL}(\g).$
\end{corollary}
\begin{proof}Again we may focus on $ \ML_{\bz}=\NL^{1+s}_{\bz} $.
We can use \cref{lem:decomposition} for an arbitrary tangent vector
$\nu\in T_\g\varOmega^{1+s}_\bz$ to obtain \cref{eq:decomposition}
with a function $\varrho\in H^{1+s}(\R/\Z)$ with $\varrho(0)=\varrho(1)=0$
and tangent vector $h\in T_\g\ML_{\bz}$. Constructing a family of
admissible parameter transformations as in the previous proof we obtain
$\D\EL_\g(\nu)=0$. Observing that we can represent 
the smooth submanifold $\varOmega_\bz^{1+s}$ as
preimage as $\varOmega_\bz^{1+s}=(\LL\times P_0)^{-1}(1,\bz)$,
where $P_0$ is the point evaluation defined in \cref{eq:point-evaluation},
we can apply the
Lagrange multiplier rule
\cite[Corollary 3.5.29]{abraham-etal_1988}
to obtain some number $\lambda\in\R$ and a vector $\Upsilon\in\R^n$ such
that $\g$ is a critical point of the functional $\EL-\lambda\LL -\langle
\Upsilon,P_0(\cdot)\rangle$. 
For the relation between $\lambda$ and $\Upsilon$ we compute similarly
as in the previous proof $ \lambda + \inner{\Upsilon,\bz}
	 =
	\lambda \, \LL(\g) + \inner{\Upsilon,P_0(\g)}
	=
	\lambda \D\LL_\g(\g) + \inner{ \Upsilon,(\D P_0)_\g(\g) }
	=
	\D\EL_\g(\g)
	= 
	\alpha \, \widetilde{\EL}(\g). $
\end{proof}

If we assume  
that the energy $\EL$ is also invariant under translation, 
we do not have to deal with the  additional
Lagrange parameter $\Upsilon\in\R^n$.

\begin{lemma}\label{cor:critical-points_wo_fixed_point}
Let $ \ML=\{\NL^{1+s},\AL^{1+s}\} $ and $ \ML_{\bz}=\{\eta\in \ML:\eta(0)=\bz\} $. If the energy $\EL\in C^1(\OL)$ is invariant under translations, then 
a critical point $\gamma \in \ML_\bz$ of $\EL|_{\ML_\bz}$ is also a 
critical point of $\EL|_{\ML}$.
\end{lemma}
\begin{proof}
	Since the energy is invariant under translations we have 
\[
	\textstyle
	\D\EL_\g p
	=
	\frac{\partial}{\partial\varepsilon}
	\EL(\gamma+\varepsilon p)
	\big|_{\varepsilon=0}
	=
	0\quad\Foa p\in\R^n
	.
\]
In particular, this implies
\[
	\textstyle
	0=\D\EL_\g h(0)
	=
	\D\EL_\g h - \D\EL_\g (h-h(0))
	=
	\D\EL_\g h
	\quad
	\Foa
	h\in T_\g\ML
	,
\]
since $h-h(0)\in T_\g\ML_\bz$;  see \cref{eq:tangent-space-AL} and
\cref{eq:tangent-space-ALz}.
\end{proof}
